\begin{solapartat}
\punts{0,2} $\vec{F} = q\vec{v}\times\vec{B}$

\punts{0,8} $\vec{F} = 1,60\cdot 10^{-19}\left(1,00\cdot 10^{5}\,\vec{\imath}\times 1,00\cdot 10^{-2}\,\vec{k}\right) = -1,60\cdot 10^{-16}\,\vec{\jmath}\ \mathrm{N}$
\end{solapartat}

\begin{solapartat}
El radi de la trajectòria del protó ve donat per:

\punts{0,4} $\left\{\begin{array}{l}
\mathit{Volem\ que}: r_p = r_e;\ qvB = m\,\dfrac{v^2}{r} \\[1ex]
v_p = v_e \\[1ex]
\dfrac{m_p v_p}{q_p B_p} = \dfrac{m_e v_e}{q_e B_e}
\end{array}\right.$

\punts{0,4} $B_e = \dfrac{m_e}{m_p}\,B_p = \dfrac{9,11\times 10^{-31}}{1,67\times 10^{-27}}\times 10^{-2} = 5,45\times 10^{-6}\ \mathrm{T}$

\punts{0,2} El camp tindrà aquest mòdul i la mateixa direcció que el primer camp però sentit contrari ja que la càrrega de l'electró té signe contrari a la del protó.

$\vec{B} = -5,45\times 10^{-6}\,\vec{k}\ \mathrm{T}$
\end{solapartat}
